\pagenumbering{roman}
\setcounter{page}{1}

\selecthungarian

%----------------------------------------------------------------------------
% Abstract in Hungarian
%----------------------------------------------------------------------------
\chapter*{Kivonat}\addcontentsline{toc}{chapter}{Kivonat}

A valós idejű deformálható testek szimulációjának néhány milliszekundumon belül
kell hihető dinamikát szolgáltatnia képkockánként, ami minden \textit{solvert} a
fizikai pontosság és a számítási költség közötti kompromisszumra kényszerít. Az
Extended Position-Based Dynamics (XPBD) évek óta a valós idejű deformálható
testek solvereinek egyik legelterjedtebb változata, míg a 2024-ben bemutatott Vertex Block Descent
(VBD) egy ígéretes alternatíva, amelyet még nem hasonlítottak össze
alaposan az XPBD-vel. Ez a szakdolgozat a két
módszer első, az eredeti cikkekhez hű, közvetlen összehasonlítását mutatja be.

Mindkét solvert ugyanabban a környezetben valósítom meg, ugyanazokra a
modellekre alkalmazom, és egyező számítási költséggel futtatom, hogy
a wall-clock költségük közvetlenül összehasonlítható legyen. Egy
magas iterációszámú Newton solver szolgál ground-truth referenciaként, így nemcsak a
relatív teljesítményről, de az abszolút
megoldási hibáról is be tudok számolni. Az
eredményeket Macklin et al. primal/dual keretrendszerén keresztül vizsgálom,
amely a két módszert ugyanazon probléma különböző irányból történő
megközelítésének tekinti, és segít az eredmények értelmezésében.

Mindkét solvert négy szempont alapján értékelem ki: hatékonyság, robusztusság,
fizikai valósághűség és rugalmasság. A fő megállapítás az, hogy egyik módszer sem
dominál; mindegyik ott teljesít jobban, ahol a frissítési iránya illik a feladathoz. A
pontosság--költség kompromisszum terén a relatív helyzetük a részlépésenkénti
feladat nehézségével változik, és mindkét módszer jellegzetes erősségeket és
hibamódokat mutat terhelés alatt. Elegendően magas iterációszámnál mindkettő eléri a
Newton-referenciát, és fizikailag megkülönböztethetetlenné válik egymástól.

A gyakorlati ajánlás tehát feladatfüggő, nem abszolút, és az eredmények
arra utalnak, hogy a többszálú vagy GPU-s végrehajtás --- amelyet az AVBD-vel
és alternatív szövetmodellekkel együtt jövőbeli munkának hagyok --- tovább
billentené a mérleget a VBD irányába.


\vfill
\selectenglish


%----------------------------------------------------------------------------
% Abstract in English
%----------------------------------------------------------------------------
\chapter*{Abstract}\addcontentsline{toc}{chapter}{Abstract}

Real-time deformable simulation has to deliver believable dynamics within a few
milliseconds per frame, which forces every solver onto a trade-off between
physical accuracy and computational cost. Extended Position-Based Dynamics
(XPBD) has been the workhorse of real-time solvers for years, while Vertex Block
Descent (VBD), introduced in 2024, is a promising alternative that has
not yet been thoroughly evaluated against it. This thesis presents the first
faithful head-to-head comparison of the two.

I implement both solvers in a shared codebase, apply them to the same models,
and run them under matched computational budgets so that their wall-clock cost
is directly comparable. A high-iteration Newton solver provides the
ground-truth reference, letting me report absolute solution error rather than
only relative performance. Throughout, I read the comparison through the
primal/dual lens of Macklin et~al., which frames the two methods as approaching
the same problem from different directions and helps interpret the results.

I evaluate both solvers along four axes: efficiency, robustness, physical
realism, and flexibility. The headline finding is that neither dominates;
each wins in the setting its update direction suits. On the accuracy--cost
trade-off their relative standing shifts with the difficulty of the per-step
problem, and each method shows characteristic strengths and failure modes
under stress. With enough iterations both converge to the Newton reference
and become physically indistinguishable.

The practical recommendation is therefore scenario-based rather than absolute,
and the results suggest that multithreaded or GPU execution --- left to future
work, together with AVBD and alternative cloth models --- would shift the
balance further toward VBD.


\vfill
\selectthesislanguage

\newcounter{romanPage}
\setcounter{romanPage}{\value{page}}
\stepcounter{romanPage}